\documentclass[11pt]{article}
\usepackage[utf8]{inputenc}
\usepackage{amsmath,amssymb,amsthm}
\usepackage{hyperref}
\usepackage{mathtools}

\title{A Three-Parameter Linear Program for the Equation $x+y=mz$}
\author{}
\date{}

\newtheorem{proposition}{Proposition}
\newtheorem{theorem}{Theorem}
\newtheorem{remark}{Remark}

\begin{document}

\maketitle

\section{Motivation}

The one-parameter and two-parameter linear programs of
\texttt{latex/3\_lp\_investigation\_notes.tex} and
\texttt{latex/4\_two\_parameter\_model.tex} plateau well above the best
published upper bounds for the density of $m$-sum-free sets.
Concretely, in all observed experiments the one-parameter LP never drops below
$1/2$, and the two-parameter LP stalls at $8/15$ for $m=3$ (see
\texttt{latex/5\_experiments\_and\_bounds.tex}).

Adding a \emph{third} parameter to the witness model enlarges the affine
symmetry group from $\mathrm{GL}_2(\mathbb{F}_p)\ltimes\mathbb{F}_p^2$ to
$\mathrm{GL}_3(\mathbb{F}_p)\ltimes\mathbb{F}_p^3$ and produces many more
useful forbidden identities. As we show below, already very small three-parameter
witness families push the LP relaxation bound strictly below $1/2$.

\section{The witness forms}

Fix a finite family of affine forms in three parameters,
\[
L_r(u,v,w)=a_r u + b_r v + c_r w + d_r,\qquad 1\le r\le n.
\]
For each sign pattern $\varepsilon\in\{\pm 1\}^n$ define the atomic set
\[
S_\varepsilon
=
\bigcap_{r=1}^n\bigl(L_r^{-1}(A)\bigr)^{\varepsilon_r}\subseteq\mathbb{F}_p^3,
\]
with $X^{+1}=X$ and $X^{-1}=\mathbb{F}_p^3\setminus X$. Its (normalized) density
$a(\varepsilon)=|S_\varepsilon|/p^3$ is the LP variable.

As usual the atoms form a probability distribution and marginalize to the
density $\delta(A)$ along any nonconstant form:
\[
\sum_\varepsilon a(\varepsilon)=1,
\qquad
a(\varepsilon)\ge 0,
\qquad
\delta(A)=\sum_{\varepsilon:\,\varepsilon_{r_0}=+1}a(\varepsilon)
\]
for any $r_0$ with $(a_{r_0},b_{r_0},c_{r_0})\ne 0$.

\section{Forbidden triples}

\begin{proposition}\label{prop:forbidden}
Assume that $(x,y,z)\in A^3\Rightarrow x+y\ne mz$. If, for a triple $(i,j,k)$,
\[
L_i(u,v,w)+L_j(u,v,w)-m L_k(u,v,w)\equiv 0
\]
as an identity in $\mathbb{F}_p^3$, then every atom with
$\varepsilon_i=\varepsilon_j=\varepsilon_k=+1$ is empty.
\end{proposition}

Coefficientwise, the hypothesis is
\[
a_i+a_j=ma_k,\quad
b_i+b_j=mb_k,\quad
c_i+c_j=mc_k,\quad
d_i+d_j=md_k.
\]
In the all-zero-shift setting $d_r=0$ this reduces to three linear equations on
the direction vectors $(a_r,b_r,c_r)\in\mathbb{Q}^3$.

\section{Affine symmetry on $\mathbb{F}_p^3$}

Any affine automorphism $T(u,v,w)=M\cdot(u,v,w)^{\mathsf T}+t$ with
$M\in\mathrm{GL}_3(\mathbb{F}_p)$ and $t\in\mathbb{F}_p^3$ preserves the
uniform distribution on $\mathbb{F}_p^3$, hence
\[
\delta\!\Bigl(\bigcap_{r\in S} L_r^{-1}(A)\Bigr)
=
\delta\!\Bigl(\bigcap_{r\in S'} L_{r'}^{-1}(A)\Bigr)
\]
whenever two subfamilies $\{L_r\}_{r\in S}$ and $\{L_{r'}\}_{r'\in S'}$ are
related by an affine reparameterization of $(u,v,w)$.

\subsection*{Rank-based normalization}

Given a subset of forms, the coefficient vectors $(a_r,b_r,c_r)$ span a linear
subspace of $\mathbb{Q}^3$ of some rank $r\in\{0,1,2,3\}$. Under the full
action of $\mathrm{GL}_3(\mathbb{F}_p)\ltimes\mathbb{F}_p^3$, we may always
rotate coordinates so that this subspace is $\mathbb{Q}^r\times\{0\}^{3-r}$.
After this, the affine model collapses to an $r$-parameter affine model, and
the canonical shape can be computed by the existing $r$-parameter logic:

\begin{itemize}
\item Rank $0$ (all forms constant): shape is the multiset of constant terms.
\item Rank $1$: project onto the 1-dimensional span and apply the one-parameter
canonical shape (\texttt{affine\_form\_lp.canonical\_subset\_shape}).
\item Rank $2$: choose any two independent forms as a basis, expand every other
form in that basis; apply the two-parameter canonical shape
(\texttt{two\_parameter\_lp.canonical\_subset\_shape}).
\item Rank $3$: try every triple of independent forms as a $\mathrm{GL}_3$ basis,
normalize to $(u,v,w)$, and take the lexicographically smallest tuple.
\end{itemize}

Tagging the canonical shape with its rank ensures that two subsets of different
ranks are never identified.

\section{The resulting LP}

Given a family $\mathcal{L}=\{L_1,\dots,L_n\}$, the three-parameter LP has

\begin{itemize}
\item variables $a(\varepsilon)\ge 0$ for $\varepsilon\in\{\pm 1\}^n$;
\item the mass constraint $\sum_\varepsilon a(\varepsilon)=1$;
\item for every identity $L_i+L_j-mL_k\equiv 0$ and every $\varepsilon$ with
$\varepsilon_i=\varepsilon_j=\varepsilon_k=+1$, the constraint
$a(\varepsilon)=0$;
\item for every pair of subsets $S,S'\subseteq\{1,\dots,n\}$ with the same
canonical shape, the equation
\[
\sum_{\varepsilon:\,\varepsilon|_S=+1}a(\varepsilon)
=
\sum_{\varepsilon:\,\varepsilon|_{S'}=+1}a(\varepsilon);
\]
\item the objective $\max\sum_{\varepsilon:\,\varepsilon_{r_0}=+1}a(\varepsilon)$.
\end{itemize}

\section{Implementation}

The module \texttt{three\_parameter\_lp.py} implements this model end-to-end:

\begin{itemize}
\item \texttt{linear\_form\_3d(a,b,c,d,label)} stores a form $au+bv+cw+d$ with
rational coefficients.
\item \texttt{canonical\_subset\_shape(forms)} computes the rank-tagged canonical
shape of a subset. Rank~1 and rank~2 subsets are reduced to the
one-parameter and two-parameter canonical shapes respectively.
\item \texttt{form\_relation\_holds(L\_i, L\_j, L\_k, m)} checks whether the
identity $L_i+L_j-mL_k\equiv 0$ holds.
\item \texttt{build\_three\_parameter\_lp(m, forms)} assembles the LP above and
solves it via SciPy's HiGHS backend.
\item \texttt{evaluate\_three\_parameter\_bound(m, forms)} returns the LP upper
bound, or \texttt{None} if the solver fails.
\end{itemize}

\section{Small illustrative families}

The following three families are defined in
\texttt{experiments\_three\_parameter.py}.

\paragraph{\texttt{basic\_pair\_family}(m).}
Six forms
\[
u,\quad v,\quad w,\quad \frac{u+v}{m},\quad \frac{u+w}{m},\quad \frac{v+w}{m}.
\]
Three forbidden triples: $u+v=m\cdot(u+v)/m$, and cyclic rotations. The LP
optimum is
\[
\boxed{\delta(A)\le \frac{16}{33}\approx 0.4848}
\]
for every $m\ge 3$ and is \emph{independent} of $m$.

\paragraph{\texttt{signed\_pair\_family}(m).}
Adds the three ``minus'' variants
\[
\frac{u-v}{m},\quad \frac{u-w}{m},\quad \frac{v-w}{m}.
\]
The same three forbidden triples hold (the new forms do not participate in any
identity $L_i+L_j=mL_k$ inside this family), but the symmetry content grows
substantially. The LP optimum is
\[
\boxed{\delta(A)\le \frac{48}{103}\approx 0.4660}
\]
uniformly in $m\ge 3$.

\paragraph{\texttt{symmetric\_12\_family}(m).}
Further adds the three ``$(2a-b-c)/m$'' forms:
\[
\frac{2u-v-w}{m},\quad \frac{-u+2v-w}{m},\quad \frac{-u-v+2w}{m}.
\]
Still only three forbidden triples, but the added symmetry drops the LP
optimum to approximately
\[
\delta(A)\le 0.4398 \;(m=3),\qquad 0.4386 \;(m=4).
\]

\section{Why the bound is uniform in $m$ for the smaller families}

The six and nine-form families use only forms of the form
$(\alpha u+\beta v+\gamma w)/m$ with $(\alpha,\beta,\gamma)\in\mathbb{Z}^3$ and
$d_r=0$. The change of coordinates $(u,v,w)\mapsto(mu,mv,mw)$ (invertible for
$\gcd(m,p)=1$) sends this family to an isomorphic family of integer forms for
every $m$. It maps forbidden identities to forbidden identities (because
$L_i+L_j=mL_k$ scales consistently), and the ranks and canonical shapes are
invariants of the affine equivalence class. Consequently the LP is the same
linear program up to a relabeling of variables, so the optimum does not depend
on $m$ for $m\ge 3$.

\section{Remarks and caveats}

\begin{remark}
These are \emph{LP relaxation} bounds for joint atomic laws compatible with
affine symmetry and the empty-atom constraints. Turning them into theorems
about $m$-sum-free sets $A\subset\mathbb{F}_p$ requires the usual step of
realizing the LP from real set statistics, as in \cite{LMPV}.
\end{remark}

\begin{remark}
With $n$ forms the LP has at most $2^n$ atomic variables; in practice the
twelve-form family already produces $2\,880$ variables after the forbidden
atoms are excluded and about $1\,700$ shape equations, and SciPy's HiGHS solves
it in a few minutes.
\end{remark}

\begin{remark}
The rank-based canonical shape is essential: the earlier three-parameter LP
stored residual $w$-coefficients for rank-2 subsets and therefore treated
``tilted'' and ``axis-aligned'' 2-planes as different orbits, losing some
symmetry equations. After projecting rank-1 and rank-2 subsets to their
natural low-dimensional models, the LP correctly identifies those subsets and
the bounds tighten from $21/41$, $11/23$ and $0.449$ down to the values listed
above.
\end{remark}

\begin{thebibliography}{1}
\bibitem{LMPV} V. Lev, M. Matolcsi, P. P. Pach, D. Varga,
\emph{On the density of Kravitz sets},
\url{https://arxiv.org/pdf/2510.16522}.
\end{thebibliography}

\end{document}
